% !TEX root = ../main.tex
\chapter{Einops: Tensor Rearrangement Algebra}
\label{ch:einops}

Einops is a small DSL for tensor rearrangement
\cite{einops_rogozhnikov}, originally developed as an external
library by Alex Rogozhnikov and adopted into the PyTorch and JAX
ecosystems. The DSL replaces several
common patterns of \code{reshape}, \code{transpose}, and
\code{view} with single readable expressions that name each axis
explicitly. \Lucid\ provides a numpy-free implementation at
\code{lucid.einops}.

This chapter describes the three operations \code{rearrange},
\code{reduce}, and \code{repeat}, the syntax of the einops
expressions, and the implementation considerations.

\section{Why Einops}
\label{sec:einops:why}

The motivating observation: code like
\code{x.transpose(0, 2).reshape(B, T * D)} hides what's happening.
The reader must work out which axes were transposed, what the
output shape means, and whether the reshape is contiguous.

The einops alternative:
\code{rearrange(x, 'b d t -> b (t d)')}.

The string is a transformation rule: the input has axes \code{b},
\code{d}, \code{t}; the output has axes \code{b} and \code{(t
d)}, which means the trailing axis is the product of \code{t}
and \code{d} (in that order). The reader sees both the input
structure and the output structure explicitly.

\subsection{The H8 namespace rule}
\label{ssec:einops:h8}

Per Hard Rule H8, einops ops are accessible only via
\code{lucid.einops.*}, never at the top level
(\code{lucid.rearrange} is invalid) and never as a Tensor method
(\code{x.rearrange(...)} is invalid). The namespace separation
matches \refframework\ and SciPy and disambiguates from the
language-level \code{reshape}.

\section{The Three Operations}
\label{sec:einops:operations}

\code{lucid.einops} exposes three operations:

\begin{itemize}
  \item \code{rearrange(tensor, pattern, **axes\_sizes)}:
    permute / reshape with an explicit axis-naming pattern.
  \item \code{reduce(tensor, pattern, reduction, **axes\_sizes)}:
    rearrange and reduce along specified axes.
  \item \code{repeat(tensor, pattern, **axes\_sizes)}: rearrange
    with axis duplication.
\end{itemize}

The \code{pattern} is always an arrow-separated string. The left
side names the input axes; the right side names the output axes.

\section{rearrange}
\label{sec:einops:rearrange}

\code{rearrange} handles pure permutations and reshapes.

\subsection{Permutation}
\label{ssec:einops:permutation}

\begin{lstlisting}[style=lucid,language=LucidPython]
from lucid.einops import rearrange

x = lucid.randn(2, 3, 4)  # shape (b, c, t)
y = rearrange(x, 'b c t -> b t c')
# y.shape == (2, 4, 3)
\end{lstlisting}

Equivalent to \code{x.permute(0, 2, 1)} but with named axes.

\subsection{Reshape via grouping}
\label{ssec:einops:reshape_group}

Parentheses group axes into a single output axis:

\begin{lstlisting}[style=lucid,language=LucidPython]
y = rearrange(x, 'b c t -> b (c t)')
# y.shape == (2, 12)   -- c and t merged
\end{lstlisting}

The order inside the parentheses matters: \code{(c t)} means
\code{c} varies more slowly than \code{t} in the merged axis,
which corresponds to a \code{reshape} of \code{(B, C, T)} to
\code{(B, C*T)} where \code{T} is the inner axis.

\subsection{Reshape via splitting}
\label{ssec:einops:reshape_split}

The inverse: split an axis using \code{axes\_sizes}:

\begin{lstlisting}[style=lucid,language=LucidPython]
y = rearrange(x, 'b (c1 c2) t -> b c1 c2 t', c1=2)
# Splits the c axis (which was 3) ... wait, this errors.
# 3 is not divisible by 2.
\end{lstlisting}

The split is only valid when the named sub-axis sizes divide the
combined axis. For \code{b (c1 c2) t} with $c1 c2 = c$, the user
must supply either \code{c1} or \code{c2} explicitly; the other
is inferred.

\subsection{Combining}
\label{ssec:einops:combining}

A single rearrange can combine permutation, splitting, and
merging:

\begin{lstlisting}[style=lucid,language=LucidPython]
# Vision Transformer patch extraction:
# (B, C, H, W) -> (B, num_patches, patch_dim)
y = rearrange(x, 'b c (h p1) (w p2) -> b (h w) (p1 p2 c)',
              p1=16, p2=16)
\end{lstlisting}

This single line replaces a \code{unfold} + \code{permute} +
\code{reshape} sequence that would be three lines and harder to
read.

\section{reduce}
\label{sec:einops:reduce}

\code{reduce} combines rearrangement with reduction:

\begin{lstlisting}[style=lucid,language=LucidPython]
from lucid.einops import reduce

# Mean-pool across H and W:
y = reduce(x, 'b c h w -> b c', 'mean')
# y.shape == (B, C), the per-channel mean

# Max-pool over patches of size 2x2:
y = reduce(x, 'b c (h p1) (w p2) -> b c h w', 'max', p1=2, p2=2)
\end{lstlisting}

The reduction is named: \code{'sum'}, \code{'mean'}, \code{'max'},
\code{'min'}, \code{'prod'}. The axes that disappear between
input and output are the reduced ones.

\subsection{vs explicit reduce}
\label{ssec:einops:vs_explicit_reduce}

The plain reduction equivalent: \code{x.mean(dim=(2, 3))} for
the first example, a more involved unfold-plus-max for the
second. \code{reduce} expresses both as one named pattern.

\section{repeat}
\label{sec:einops:repeat}

\code{repeat} expands axes with named multiplication:

\begin{lstlisting}[style=lucid,language=LucidPython]
from lucid.einops import repeat

# Add a batch dim of size B by repeating:
x = lucid.randn(C, H, W)
y = repeat(x, 'c h w -> b c h w', b=8)
# y.shape == (8, C, H, W), with the c h w part identical across b
\end{lstlisting}

The \code{b} axis is new (not in the input); the size is given by
the keyword argument. The data is virtually duplicated via
stride-0 broadcasting; the user can pass through to a kernel that
respects strides, or call \code{.contiguous()} to materialise.

\section{Pattern Syntax}
\label{sec:einops:syntax}

The pattern grammar is small but expressive.

\subsection{Axis names}
\label{ssec:einops:axis_names}

An axis name is an identifier: letters, digits, underscores;
starts with a letter. Common conventions:

\begin{itemize}
  \item \code{b}: batch.
  \item \code{c}: channels.
  \item \code{h}, \code{w}: height, width.
  \item \code{t}: time.
  \item \code{n}: a generic count.
\end{itemize}

The names have no semantic meaning to einops; they are just
labels. The convention is mnemonic for the reader.

\subsection{Grouping}
\label{ssec:einops:grouping}

\code{(a b)} means a single axis whose size is the product of
\code{a} and \code{b}. The order matters: \code{a} is the slow
axis, \code{b} is the fast.

\subsection{Ellipsis}
\label{ssec:einops:ellipsis}

\code{...} matches a variable number of axes:

\begin{lstlisting}[style=lucid,language=LucidPython]
# Move the last axis to the front, regardless of rank:
y = rearrange(x, '... c -> c ...')
\end{lstlisting}

\subsection{The 1 placeholder}
\label{ssec:einops:one_placeholder}

The literal \code{1} represents a size-1 axis:

\begin{lstlisting}[style=lucid,language=LucidPython]
# Unsqueeze the second axis:
y = rearrange(x, 'b t -> b 1 t')
\end{lstlisting}

\section{Backward Through Einops}
\label{sec:einops:backward}

All three einops operations are differentiable. The backward is
the inverse pattern:

\begin{itemize}
  \item \code{rearrange}'s backward applies the inverse
    permutation/reshape.
  \item \code{reduce}'s backward broadcasts the gradient back
    along the reduced axes (with the appropriate scaling for
    \code{'mean'}).
  \item \code{repeat}'s backward sums the gradient over the
    repeated axes.
\end{itemize}

The implementation composes existing \Lucid\ primitives, so the
autograd graph is the same as if the user had written the
equivalent permute / sum / etc.\ sequence by hand.

\section{Implementation}
\label{sec:einops:impl}

The implementation lives in \code{lucid/einops/} as pure Python.
The parser converts the pattern string to a sequence of
\code{reshape} / \code{permute} / \code{reduce} / \code{expand}
operations on the input tensor.

\subsection{Parser}
\label{ssec:einops:parser}

The parser tokenises the pattern into axis names and
parenthesised groups, validates the input and output sides for
consistency (every input axis appears somewhere in the output, or
is reduced), and produces a sequence of tensor ops.

The parser is small (a few hundred lines) and produces clear
error messages for malformed patterns.

\subsection{Caching}
\label{ssec:einops:caching}

For a given (pattern, input shape, axes\_sizes) tuple, the
sequence of ops is deterministic. \Lucid\ caches the parsed ops
by these keys so that repeated calls in a training loop pay the
parse cost only once.

The cache is a small dict keyed by hash of the inputs; size is
bounded by an LRU policy (default 1024 entries, which covers any
reasonable model).

\section{Patterns in Practice}
\label{sec:einops:patterns_in_practice}

A few patterns that recur in ML.

\subsection{Attention reshape}
\label{ssec:einops:attention_reshape}

Multi-head attention partitions the channel dimension into
\code{h} heads of \code{d} dimensions each:

\begin{lstlisting}[style=lucid,language=LucidPython]
# (B, T, H*D) -> (B, H, T, D)
qkv = rearrange(qkv, 'b t (h d) -> b h t d', h=n_heads)
\end{lstlisting}

\subsection{Patch embedding (ViT)}
\label{ssec:einops:patch_embed}

Vision Transformer's patch extraction:

\begin{lstlisting}[style=lucid,language=LucidPython]
# (B, C, H, W) -> (B, (H/p)(W/p), p*p*C)
patches = rearrange(image, 'b c (h p1) (w p2) -> b (h w) (p1 p2 c)',
                    p1=16, p2=16)
\end{lstlisting}

\subsection{Channel-last conversion}
\label{ssec:einops:nhwc}

\begin{lstlisting}[style=lucid,language=LucidPython]
# NCHW (PyTorch default) <-> NHWC (TF default)
y = rearrange(x, 'b c h w -> b h w c')
\end{lstlisting}

\section{Comparison: Einops vs Native Operations}
\label{sec:einops:vs_native}

To make the readability claim concrete, several common operations
expressed both ways.

\subsection{Multi-head attention reshape}
\label{ssec:einops:mha_compare}

Going from $(B, T, H \cdot D)$ to $(B, H, T, D)$ for multi-head
attention:

\begin{lstlisting}[style=lucid,language=LucidPython]
# Native
qkv = qkv.view(B, T, H, D).permute(0, 2, 1, 3)

# Einops
qkv = rearrange(qkv, 'b t (h d) -> b h t d', h=H)
\end{lstlisting}

The native version requires the reader to recall that
\code{view(B, T, H, D)} splits the last dim into $H$ and $D$ (and
that $H$ comes before $D$ in the split), then permutes to put the
heads in dim 1. The einops version makes both transformations
visible.

\subsection{Patch extraction (ViT)}
\label{ssec:einops:patch_compare}

\begin{lstlisting}[style=lucid,language=LucidPython]
# Native (using unfold)
patches = (image
    .unfold(2, P, P).unfold(3, P, P)
    .contiguous()
    .view(B, C, n_patches_h, n_patches_w, P, P)
    .permute(0, 2, 3, 4, 5, 1)
    .reshape(B, n_patches_h * n_patches_w, P * P * C))

# Einops
patches = rearrange(image, 'b c (h p1) (w p2) -> b (h w) (p1 p2 c)',
                    p1=P, p2=P)
\end{lstlisting}

The native version is six lines, four operations, and requires
the reader to follow the dimension order through each step. The
einops version is one line.

\subsection{Channel-last conversion with batch}
\label{ssec:einops:nhwc_compare}

\begin{lstlisting}[style=lucid,language=LucidPython]
# Native
x = x.permute(0, 2, 3, 1).contiguous()

# Einops
x = rearrange(x, 'b c h w -> b h w c')
\end{lstlisting}

Native is shorter here, but the einops version makes the
target layout explicit; the reader does not have to count the
permutation indices.

\section{Pattern Parsing in Detail}
\label{sec:einops:parsing}

The parser converts a pattern string into a sequence of tensor
operations. This section sketches the algorithm.

\subsection{Tokenisation}
\label{ssec:einops:tokenisation}

The pattern is split on whitespace and parenthesisation. Each
token is one of:

\begin{itemize}
  \item An axis name (identifier).
  \item A group: parenthesised sequence of axis names.
  \item The ellipsis \code{...}.
  \item The literal \code{1}.
\end{itemize}

The input side and output side are separated by \code{->}.

\subsection{Axis-size resolution}
\label{ssec:einops:size_resolution}

Each named axis must have a determined size at evaluation time.
Sizes come from three sources:

\begin{itemize}
  \item The input tensor's shape: an axis name appearing alone
    on the input side gets its size from the corresponding input
    dimension.
  \item Keyword arguments: \code{rearrange(x, '...', h=4)} sets
    the axis named \code{h} to size 4.
  \item Inference: in a group, one axis can be inferred from the
    others. \code{(h d)} for an input dim of size 12 with $h = 4$
    given gives $d = 3$.
\end{itemize}

A pattern with unresolved axes raises a parse error.

\subsection{Op sequence generation}
\label{ssec:einops:op_sequence}

The parser produces a sequence of tensor operations to transform
the input into the output:

\begin{enumerate}
  \item \emph{Split} input groups: for each input \code{(a b)},
    apply \code{reshape} that turns the combined dim into two.
  \item \emph{Permute}: reorder the dimensions so they appear in
    the output order.
  \item \emph{Reduce} (for \code{reduce}): apply the named
    reduction along the axes that disappear.
  \item \emph{Expand} (for \code{repeat}): apply \code{expand} or
    \code{repeat} to introduce new axes.
  \item \emph{Merge} output groups: for each output \code{(a b)},
    apply \code{reshape} that flattens.
\end{enumerate}

The op sequence is cached by (pattern, input rank, axis sizes)
tuple so that repeated calls in a training loop skip the parse.

\section{Performance Considerations}
\label{sec:einops:performance}

\subsection{Overhead per call}
\label{ssec:einops:overhead_per_call}

The first call with a new (pattern, shape) tuple pays the parse
cost (a few hundred microseconds). Subsequent calls hit the cache
in a few microseconds.

For training loops where the same patterns are called thousands
of times per step, the cached path is fast enough to be
negligible relative to the underlying kernel cost. For very
small kernels (sub-microsecond), the einops overhead can be
noticeable; for any non-trivial kernel, it is dominated.

\subsection{When einops materialises}
\label{ssec:einops:when_materialise}

The op sequence is a chain of \code{reshape}, \code{permute},
\code{expand}, etc. Each is a view operation when possible. The
only operation that always allocates is \code{repeat} (physical
duplication), which the user can avoid by using \code{expand} +
care.

\subsection{Compose with autograd}
\label{ssec:einops:compose_autograd}

Each underlying primitive is differentiable; the autograd graph
is the same as if the user had written the equivalent native
sequence. There is no extra cost or memory overhead beyond what
the native sequence would incur.

\section{Advanced Patterns}
\label{sec:einops:advanced_patterns}

\subsection{Pattern with ellipsis}
\label{ssec:einops:ellipsis_pattern}

\code{...} matches zero or more leading dimensions:

\begin{lstlisting}[style=lucid,language=LucidPython]
# Sum the last axis regardless of leading rank
y = reduce(x, '... c -> ...', 'sum')
\end{lstlisting}

\subsection{Pattern with literal 1}
\label{ssec:einops:literal_one}

The literal \code{1} represents a size-1 axis. Useful for
explicit broadcasting:

\begin{lstlisting}[style=lucid,language=LucidPython]
# Insert a size-1 axis for broadcasting
expanded = rearrange(x, 'b c -> b c 1')
result = expanded + other
\end{lstlisting}

\subsection{Pattern with anonymous axis}
\label{ssec:einops:anon_axis}

An underscore prefix marks an axis as anonymous (matches but is
not constrained to a particular name):

\begin{lstlisting}[style=lucid,language=LucidPython]
# Move the last anonymous axis to first
y = rearrange(x, '... _final -> _final ...')
\end{lstlisting}

In practice, named axes are clearer; the anonymous form is
rare.

\section{The Three Operations: A Reference}
\label{sec:einops:reference}

\tabref{tab:einops:reference} is the complete reference for
\code{rearrange}, \code{reduce}, and \code{repeat}.

\begin{table}[t]
\centering
\small
\begin{tabular}{lll}
\toprule
Operation & Effect on shape & Use case \\
\midrule
\code{rearrange} & no axis added or removed & permute, reshape, group, split \\
\code{reduce}   & axes disappear (those on input only) & per-axis aggregation \\
\code{repeat}   & axes appear (those on output only) & broadcast-like duplication \\
\bottomrule
\end{tabular}
\caption[Einops operation reference.]{The three einops operations
and their effect on the tensor shape. The rule of thumb: an axis
present on input only is reduced; present on output only is
repeated; present on both (or in a group) is preserved.}
\label{tab:einops:reference}
\end{table}

\section{Common Bugs and Their Diagnostics}
\label{sec:einops:bugs}

\subsection{Bug: unresolvable size}
\label{ssec:einops:bug_size}

\code{rearrange(x, 'b (h d) -> b h d')} for $x$ of shape
$(2, 12)$ requires one of $h$ or $d$ to be supplied. Without
that, the parser raises \code{ParseError: cannot infer size of h
or d}.

Fix: \code{rearrange(x, 'b (h d) -> b h d', h=4)}.

\subsection{Bug: mismatched output axes}
\label{ssec:einops:bug_mismatch}

\code{rearrange(x, 'b c h w -> b c h')} drops $w$ on the output
side without a reduction. The parser raises
\code{ParseError: axis 'w' appears on input but not on output (use reduce?)}.

Fix: use \code{reduce(x, 'b c h w -> b c h', 'mean')} or
similar.

\subsection{Bug: division-incompatible split}
\label{ssec:einops:bug_div}

\code{rearrange(x, 'b (h d) -> b h d', h=3)} for $x$ of shape
$(2, 7)$: the dimension of size 7 cannot be split into 3 by $d$.
The parser raises \code{ShapeError: 7 is not divisible by 3}.

\section{Summary and Forward References}
\label{sec:einops:summary}

\code{lucid.einops} provides three operations
(\code{rearrange}, \code{reduce}, \code{repeat}) with a string-
based pattern syntax that names axes explicitly. The result is
more readable than the equivalent \code{reshape} /
\code{permute} / \code{sum} sequence; the implementation
composes existing primitives and is automatically differentiable.

This chapter closes \partref{part:math_subsystems}.
\partref{part:neural_networks} begins the neural-network library:
modules, layers, attention, losses, optimisers, and AMP.

\begin{lucidnote}
For the user: einops is a convenience layer. Anything it does
could be written with \code{reshape}, \code{permute}, and
\code{sum} directly. Einops just makes the intent visible. For
attention reshapes, patch embeddings, and channel-axis
conversions, the readability win is substantial.
\end{lucidnote}
